% ===========================================================================
% chapters/chap01_intro.tex — 第一章 绪论（示例章节）
% 本文件展示：
%   - 正文段落样式（宋体小四、首行缩进2字符、行距18磅）
%   - 各级标题样式
%   - 图表插入示例（含标题格式）
%   - 公式编号示例
%   - 交叉引用示例
% ===========================================================================

\chapter{绪论}
\label{chap:intro}

% ===== 二级标题 =====
\section{研究背景与意义}
\label{sec:background}

% ===== 三级标题 =====
\subsection{研究背景}
\label{subsec:bg}

粮食产量的变化对全球社会的稳定和人民的基本需求有着重大影响。《2025 年世界粮食安全和营养状况》中提到，相较于 2023 年和 2022 年，2024 年全球饥饿人口规模呈稳步下降趋势\cite{ref1}，但数量仍在 6.38 亿至 7.2 亿之间，占比约 8.2\%。

我国是世界上的农业大国，始终将"三农"工作作为重中之重。2026 年《中共中央国务院关于锚定农业农村现代化扎实推进乡村全面振兴的意见》中明确提出了"稳定发展粮油生产""提升农业科技创新效能"等重要要求。由此可见，国家对抓好粮食和重要农产品生产、加强农业基础设施建设和保障国家粮食安全等方面高度重视。

水稻作为我国最主要的粮食作物之一，其产量和品质直接关系到国家粮食安全和人民生活水平。精准监测水稻物候信息并及时估测产量，对于制定科学的田间管理措施、优化农业资源配置、保障粮食供给具有重要的现实意义。

% --- 图片插入示例 ---
% 对应格式要求：图序及图名置于图的下方，居中
% 标题格式：小五号黑体，中英文对照
% 图与正文之间要有一行的间距
\begin{figure}[htbp]
  \centering
  % 请将示例图片放在项目根目录下
  % 可使用 example-image（LaTeX 自带示例图片）测试
  \includegraphics[width=0.7\textwidth]{example-image}
  \caption{2019--2023年我国水稻播种面积、产量和单位产量}
  \label{fig:rice_stats}
  % 英文标题
  \bigskip
  \centering
  {\zihao{5}\normalfont Figure~\ref{fig:rice_stats} The sown area, yield, and unit yield of rice in China from 2019 to 2023}
\end{figure}

\subsection{研究意义}
\label{subsec:significance}

目前仅依托图像识别技术开展物候期判别或产量估测，易受拍摄角度、背景干扰等因素影响，导致数据质量不稳定，无法充分满足实际研究与应用需求，尤其是无法满足早期估产的需求。光谱成像技术能够通过多波段信息捕捉更为全面地反映水稻冠层信息，为物候监测和产量估测提供更丰富、更精准的数据支撑。

无人机低空遥感成本低、数据获取周期可控，且能通过搭载不同类型的传感器得到多源多模态数据。因此，使用无人机遥感提供的光谱信息研究物候期的判别与水稻产量估测具有重要意义。

% ===== 二级标题 =====
\section{国内外研究现状分析及存在问题}
\label{sec:literature}

\subsection{作物物候监测识别现状分析}
\label{subsec:phenology}

作物物候期为农田管理人员开展灌溉、施肥和收获等农田作业提供重要参考。准确的物候信息一方面能够实现高效的田间管理，进而提升作物产量，另一方面可为选择精准的产量估测模型提供依据\cite{ref2}。目前作物物候监测的方式主要包括人工抽样调查、基于计算机视觉技术的图像监测识别、基于无人机遥感和卫星遥感等平台的光学反演等。

% --- 表格插入示例 ---
% 对应格式要求：表序及表名置于表的上方，居中
% 表格一律用三线表（booktabs 宏包）
% 标题格式：小五号黑体，中英文对照
\begin{table}[htbp]
  \centering
  \caption{近年作物物候监测识别研究情况}
  \label{tab:phenology_studies}
  \bigskip
  {\zihao{5}\normalfont Table~\ref{tab:phenology_studies} Recent studies on crop phenology monitoring and identification}\\[4pt]
  \begin{tabular}{lccc}
    \toprule
    研究 & 作物 & 数据源 & 最优方法 \\
    \midrule
    Feng~等人 & 玉米 & 无人机 & 最小二乘法 \\
    Zhu~等人 & 玉米 & MODIS & 双特征点加权法 \\
    Zhang~等人 & 水稻 & 无人机 & 随机森林 \\
    Liu~等人 & 小麦 & Sentinel-2 & 深度学习 \\
    \bottomrule
  \end{tabular}
\end{table}

\subsection{作物产量估测现状分析}
\label{subsec:yield}

作物产量的准确估测对于农业政策制定、粮食市场调控和农业生产管理具有重要意义。传统的产量估测方法主要基于统计模型、作物生长模型以及卫星遥感数据等方法。近年来，随着机器学习和深度学习技术的快速发展，基于多源遥感数据融合的产量估测方法成为研究热点。

% --- 公式示例 ---
% 对应格式要求：公式居中，编号用括号括起写在右边行末
% 公式与正文之间要有一行的间距
随机森林回归模型可表示为：
\begin{equation}
  \label{eq:rf}
  \hat{y} = \frac{1}{N} \sum_{i=1}^{N} T_i(\mathbf{x})
\end{equation}
其中，$\hat{y}$ 为预测产量值，$N$ 为决策树数量，$T_i(\mathbf{x})$ 为第 $i$ 棵决策树的预测结果。

深度神经网络模型的训练目标是最小化损失函数：
\begin{equation}
  \label{eq:loss}
  \mathcal{L}(\theta) = \frac{1}{M} \sum_{j=1}^{M} \left( y_j - f(\mathbf{x}_j; \theta) \right)^2
\end{equation}
其中 $\theta$ 为网络参数，$M$ 为样本数量。

如公式~\eqref{eq:rf} 和公式~\eqref{eq:loss} 所示，集成学习方法和深度学习方法在产量估测方面各有优势。

\section{本章小结}
\label{sec:summary}

本章阐述了水稻物候监测与产量估测的研究背景和意义，从物候监测、产量估测两个方面系统综述了国内外研究现状，分析了现有研究存在的问题和技术难点，明确了本研究的主要内容和目标。
